\documentclass[10pt,a4paper]{amsart}
\usepackage[margin=19mm]{geometry}
\usepackage{amsmath,amssymb,mathtools}
\usepackage[scr=rsfs,cal=euler]{mathalfa}
\usepackage{xfrac}
\usepackage[unicode,hidelinks,bookmarks=false]{hyperref}
\newcommand{\ie}{i.e.\ }
\newcommand{\cf}{cf.\ }
\newcommand{\resp}{respectively\ }
\newcommand{\Subsection}{\subsection}
\providecommand{\nobreakdash}{\nobreak\hbox{-}}
\setlength{\emergencystretch}{2em}
\multlinegap0pt
\usepackage{csquotes}
\usepackage{amssymb}
\usepackage{enumerate}
\usepackage{mathtools}
\usepackage[shortlabels]{enumitem}
\usepackage{cancel}
\usepackage{tikz}
\usepackage{dsfont}
\usepackage[normalem]{ulem}
\usepackage{bm}
\usepackage{tcolorbox}
\newtheorem{proposition}{Proposition}
\usepackage{cleveref}
\crefname{proposition}{Proposition}{Propositions}
\newcommand{\CV}{\mathbf{C}_V}
\newcommand{\CVp}[1]{\mathbf{C}_{V,#1}}
\newcommand{\CWp}[1]{\mathbf{C}_{W,#1}}
\newcommand{\ThetaPar}{\Theta}
\newcommand{\cVp}[1]{\mathbf{c}_{V,#1}}
\newcommand{\cW}{\mathbf{c}_W}
\renewcommand{\epsilon}{\varepsilon}
\title{A dimension-free interpolation\\ of Caffarelli’s contraction theorem}
\author{Bader Ammari, Alessio Figalli}
\date{Source: arXiv:2605.24443v1 (CC BY 4.0)}
\begin{document}
\maketitle

\noindent\emph{Source and licence.} A dimension-free interpolation\\ of Caffarelli’s contraction theorem. Version arXiv:2605.24443v1 (CC BY 4.0), distributed by arXiv under the Creative Commons Attribution 4.0 International licence, http://creativecommons.org/licenses/by/4.0/, as recorded in the arXiv licence metadata for this exact version and shown on the abstract page https://arxiv.org/abs/2605.24443v1. Full licence notice: \texttt{licenses/arxiv-2605.24443-CC-BY-4.0.txt}.\par\smallskip

\noindent\textit{Editorial scope.} Counted unit: the article's proposition Conditional\_finite\_D (source0.tex lines 900--942). The article's complete original proof of it is delivered with it (source0.tex lines 944--1021, one \texttt{proof} environment of 78 lines). No mathematics is written, repaired or re-derived: the statement and the proof are the article's own lines, in the article's own notation, and no retained line is substituted or re-flowed.  Retained uncounted as the source-local context the proof needs: lines 452--461; lines 615--623; lines 627--636.  Nothing was omitted from any retained span, so the package carries no omission record.  No retained line contains a citation, so the wrapper carries no bibliography and no bibliography entry.  No retained line contains a figure environment, an image-inclusion command or drawing code, so no image is delivered and the package declares no asset dependency.  Editorial formatting only, in addition: an AMS \texttt{amsart} preamble that restates the article's own declarations and macros used by the retained lines, namely the environments \texttt{proposition} on separate counters, together with the standard packages those lines need.  The retained environments print in this document as Proposition 1 in order of appearance, whereas the article prints the same items with the numbering of its own full document.  The complete native source of the article as distributed in the arXiv e-print for this exact version is bundled unchanged as \texttt{sources/201274/source0.tex}.
\begin{equation}\label{Gamma_definition}
\Gamma_{d,D}:=
\begin{cases}
\displaystyle
+\infty, & D=d<+\infty,\\[0.8em]
\displaystyle
\left(\frac{2d-D}{D-d}\right)_+, & d<D<+\infty,\\[0.8em]
0, & D=+\infty.
\end{cases}
\end{equation}

\begin{equation}\label{Exact_second_variation}
\begin{aligned}
0 \ge{}&
\ThetaPar_D'(W)W_{11}\lambda^2-\ThetaPar_d'(V)V_{11}
-2\ThetaPar_d''(V)V_1^2 \\
&+\left(\frac{\ThetaPar_D'(W)^2}{d}+\ThetaPar_D''(W)\right)W_1^2\lambda^2
-\frac{2\ThetaPar_d'(V)\ThetaPar_D'(W)}{d}V_1W_1\lambda,
\end{aligned}
\end{equation}

\begin{equation}\label{Master_second_variation}
    \begin{aligned}
    W_{11}\lambda^2
    \le{}&
    \frac{1}{1-\epsilon}\frac{d}{D}\frac{D+W}{d+V}V_{11} \\
    &+
    \frac{1}{\epsilon(1-\epsilon)}
    \frac{V_1^2W_1^2}{W_{11}(d+V)^2},
    \end{aligned}
\end{equation}

\begin{proposition}\label{Conditional_finite_D}
Let $n\le d\le D<+\infty$ and $0<R\le+\infty$. Assume that $V>-d$, $W>-D$, and
\[
    D^2V\le \CV^{(2)}\,\mathrm{Id}, \qquad
    D^2W\ge \cW^{(2)}\,\mathrm{Id},
\]
with $0<\cW^{(2)},\CV^{(2)}<+\infty$. Let $\nabla\varphi$ be the optimal transport map pushing $\rho_{V,d}$ forward to $\rho_{W,D}$. Assume that, for some $\mathsf R_R,\mathsf G_R\in[0,+\infty)$,
\begin{equation}\label{Growth_assumption_general}
    |\nabla\varphi(x)|\le \mathsf R_R,
    \qquad
    \frac{d}{D}\frac{D+|\nabla\varphi(x)|^2}{d+|x|^2}\le \mathsf G_R
    \qquad \text{for a.e. }\, x\in B_R(0).
\end{equation}
Then
\begin{equation}\label{Finite_D_local_bound}
    \|D^2\varphi\|_{L^\infty(B_R(0))}
    \le
    \sqrt{\widetilde A_{V,W}^{d,D}(R)+\widetilde B_{V,W}^{d,D}(R)}+\sqrt{\widetilde B_{V,W}^{d,D}(R)},
\end{equation}
where, with $\Gamma_{d,D}$ as in \cref{Gamma_definition}, we set
\[
\widetilde\Lambda_R
:=
\CWp{D}^{(0)}(\mathsf R_R)\mathsf G_R,\qquad
\widetilde\Xi_R
:=
\min\left\{
\frac{D}{d}\,
\frac{\CWp{D}^{(1)}(\mathsf R_R)}
{\CWp{D}^{(0)}(\mathsf R_R)},
\cW^{(2)}\Gamma_{d,D}
\right\}
\]
\[
\widetilde A_{V,W}^{d,D}(R)
=
    \frac{\CV^{(2)}\widetilde\Lambda_R}{\cW^{(2)}\cVp{d}^{(0)}(R)},\qquad
\widetilde B_{V,W}^{d,D}(R)
=
    \frac{4\CVp{d}^{(1)}\widetilde\Lambda_R\widetilde\Xi_R}
    {(\cW^{(2)})^2(\cVp{d}^{(0)}(R))^2}.
\]
\end{proposition}

\begin{proof}
As discussed above, we prove the result under the assumption that there exists an interior  maximum point $(\overline x,e_1)$ with $\overline x\in B_R(0)$. This assumption can be removed by using the localized incremental-ratio argument described at the end of the paper.

Set $\overline y=\nabla\varphi(\overline x)$. By \cref{Master_second_variation} and the Hessian bounds,
\[
\cW^{(2)}\varphi_{11}(\overline x)^2
\le
\frac{\CV^{(2)}}{1-\epsilon}
\frac{d(D+W(\overline y))}{D(d+V(\overline x))} 
+
\frac{1}{\epsilon(1-\epsilon)\cW^{(2)}}
\frac{V_1(\overline x)^2W_1(\overline y)^2}{(d+V(\overline x))^2}.
\]
Using the definitions of the constants and the growth assumption \cref{Growth_assumption_general}, we have
\[
    \frac{d}{D}\frac{D+W(\overline y)}{d+V(\overline x)}
    \le
    \frac{\CWp{D}^{(0)}(\mathsf R_R)\mathsf G_R}{\cVp{d}^{(0)}(R)}
\]
and
\[
    \frac{V_1(\overline x)^2W_1(\overline y)^2}{(d+V(\overline x))^2}
    \le
    \frac{4\CVp{d}^{(1)}\CWp{D}^{(1)}(\mathsf R_R)\frac{D}{d}\mathsf G_R}{(\cVp{d}^{(0)}(R))^2}.
\]
Therefore, with
\[
    B_y:=
    \frac{4\CVp{d}^{(1)}\CWp{D}^{(1)}(\mathsf R_R)\frac{D}{d}\mathsf G_R}
    {(\cW^{(2)})^2(\cVp{d}^{(0)}(R))^2},
\]
we have
\[
    \begin{aligned}
    \varphi_{11}(\overline x)^2
    \le{}&
    \frac{\widetilde A_{V,W}^{d,D}(R)}{1-\epsilon}
    +
    \frac{B_y}{\epsilon(1-\epsilon)}.
    \end{aligned}
\]
Optimizing in $\epsilon\in(0,1)$ gives the desired estimate with $B_y$ in place of $\widetilde B_{V,W}^{d,D}(R)$.

If $D>d$, we can instead use the positive terms $W_1^2\lambda^2$ and $V_1^2$ in \cref{Exact_second_variation}. Minimizing the resulting quadratic polynomial in $W_1\lambda$ gives
\[
    \ThetaPar_D'(W)W_{11}\lambda^2
    \le
    \ThetaPar_d'(V)V_{11}
    +
    \frac{d}{(d+V)^2}
    \left(\frac{2d-D}{D-d}\right)_+
    V_1^2.
\]





Using again the definitions of the constants and \cref{Growth_assumption_general}, this yields
\[
    \varphi_{11}(\overline x)^2
    \le
    \widetilde A_{V,W}^{d,D}(R)
    +
    B_c,
\]
where
\[
    B_c:=
    \frac{4\CVp{d}^{(1)}\CWp{D}^{(0)}(\mathsf R_R)\mathsf G_R\Gamma_{d,D}}
    {\cW^{(2)}(\cVp{d}^{(0)}(R))^2}.
\]
For $D=d$, we use the convention $B_c=+\infty$. Since the function
\[
    B\mapsto \sqrt{\widetilde A_{V,W}^{d,D}(R)+B}+\sqrt B
\]
is increasing, the two bounds give \cref{Finite_D_local_bound} with $B=\min\{B_y,B_c\}$, namely with $\widetilde B_{V,W}^{d,D}(R)$.
\end{proof}

\end{document}
